\documentclass[12]{article}
\usepackage{tikz}
\usetikzlibrary{positioning}
\definecolor{processblue}{cmyk}{0.96, 0, 0, 0}
\usepackage{import}
\usepackage{transparent}
\usepackage{xcolor}
\usepackage{tcolorbox}
\usepackage{amsmath}
\usepackage{upgreek}
\usepackage{enumitem}
\usepackage{amssymb}
\usepackage[margin=0.25in]{geometry}
\usepackage{listings}
\usepackage[T1]{fontenc}
%\geometry{letterpaper}
%Commands definitions
\newcommand{\setbackgroundcolour}{\pagecolor[rgb]{0.19, 0.19, 0.19}}
\newcommand{\settextcolour}{\color[rgb]{0.87, 0.87, 0.87}}
\newcommand{\invertbackgroundtext}{\setbackgroundcolour\settextcolour}

%Command execution
%If this line is commented, then the appearance remaind as usual.
\invertbackgroundtext

\lstset{
    language=C,
    basicstyle=\ttfamily,
    keywordstyle=\color{cyan},
    commentstyle=\color{green!70!black},
    stringstyle=\color{red},
    showstringspaces=false,
    tabsize=4,
    breaklines=true,
    breakatwhitespace=false,
}

% Ponizsze linijki sa templatka do zadan i problemow                                                                                                           
% Wziete gdzies z internetu
% Assignment class for homework and exams.                                                                                                                  
%
% Gabriel Antonius
%
%
\NeedsTeXFormat{LaTeX2e}
\RequirePackage{placeins}
\RequirePackage{environ}
\RequirePackage{xifthen}


% \solutiontrue and \solutionfalse control whether solutions are shown
\newif\ifsolution
\solutiontrue
%\solutionfalse


% Redefine these if you want to change the language
\newcommand{\problemlabel}{Zadanie}
\newcommand{\solutionlabel}{Rozwi\k{a}zanie}


% Set up counters for problems and subproblems
\newcounter{ProblemNum}
\newcounter{SubProblemNum}[ProblemNum]
\renewcommand{\theProblemNum}{\arabic{ProblemNum}}
\renewcommand{\theSubProblemNum}{\alph{SubProblemNum}}


% The problem environment is the base unit of content for this class.
\newcommand{\subsectiontitle}{}
\newenvironment{problem}[1]%
  {
  \stepcounter{ProblemNum}
  \renewcommand{\subsectiontitle}{\problemlabel \ \theProblemNum \ifthenelse{\isempty{#1}}{}{: #1}}
  \medskip \subsection*{\subsectiontitle}  % new line after title
  {\large Tre\'{s}\'{c}: } 
  \FloatBarrier
  }
  {
  \FloatBarrier
  }


% The subproblem command divides a problem into parts a), b), c), ...
\newcommand*{\subproblem}{%
  \vspace*{3pt}
  \stepcounter{SubProblemNum}%
  {\bf \theSubProblemNum)\hspace{2pt}}
  }


% The solution environment should be used within the problem environment
\NewEnviron{solution}{
  \setcounter{SubProblemNum}{0}
  \ifsolution
    \FloatBarrier
    \subsubsection*{\solutionlabel}
    \BODY
  \fi}


%\documentclass[12pt]{myassignment}
%\usepackage{amsmath}                                                                                                                                        
\setlength{\parindent}{0pt}

\solutiontrue  % <---- Toggle showing solutions or not
%\solutionfalse

%koniec templatki do zadan i problemow




%\author{Lukasz Kopyto}
%\title{tytul}

\begin{document}
%\maketitle
\noindent \L ukasz Kopyto \hfill 12.12.2024 \\

\begin{center}

    {\large Notki do listy 5}
\end{center}
\section*{Twierdzenie Eulera i pewne jego zastosowania}


Niech $m > 1$ i $m\in \mathbb{N}$ oraz niech

\[r_{1}, r_{2}, \dots, r_{\phi(m)} \tag{1}\]

bedzie ciagiem wszystkich liczb naturalnych $\leq m$ i wzglednie pierwszych do $m$. Niech a bedzie jakakolwiek liczba calkowita wzglednie pierwsza do $m$. Dla $k=1,2,\dots, \phi(m)$ oznaczmy przez $\delta_{k}$ reszte z dzielenia liczby $ar_{k}$ przez $m$. Bedzie wiec:

\[ \delta_{k}\equiv ar_{k}\text{(mod m) dla } k = 1,2,\dots \phi(m)\tag{2}\]

Pokazemy, ze ciag (2) rozni sie co najwyzej porzadkiem wyrazow do ciagu (1). W tym celu wystarczy pokazac, ze
\begin{enumerate}
    \item kazdy wyraz ciagu (2) jest liczba naturalna $\leq m$ i wzglednie pierwsza do m
    \item wszystkie wyrazy ciagu (2) sa rozne
\end{enumerate}

Poniewaz liczby z ciagu (2) sa resztami z dzielenie przez m, wiec oczywiscie sa one $\geq 0$ oraz $< m$. Wobec (2) mamy $\delta_{k} = ar_{k} + mt_{k}$, dla $k = 1, 2, \dots, \phi(m)$, gdzie $t_{k}$ jest liczba calkowita. Wobec $gcd(a, m) = 1$ oraz $gcd(r_{k}, m) = 1$ dla $k = 1,2,\dots, \phi(m)$ mamy $gcd(ar_{k}, m) = 1$. Gdyby bylo tak, ze $d | \delta_{k}$ i $d|m$ to wobec wzoru na $\delta_{k}$ mielibysmy $d|ar_{k}$, co wobec $d|m$ i $gcd(ar_{k}, m) = 1$, daje $d|1$. Jest wiec $gcd(\delta_{k}, m)= 1$. Warunek 1. jest wiec spelniony.

Przypuscmy teraz, ze przy pewnych liczbach k i l, gdzie $1\leq k < l < \phi(m)$, mamy $\delta{k} = \delta_{1}$. Wobec (2) byloby $ar_{k} \equiv ar_{1}\text{(mod m)}$, skad $m| (r_{1} - r_{k})$ co, wobec $gcd(a,m) = 1$ daje $m| r_{1} - r_{k}$.

Ale $r_{k}$ i $r_{1}$ sa to liczby naturalne < m, zatem wartosc bezwzgledna ich roznicy jest < m. i z ostatniego wzoru wynika, ze $r_{k} = r_{l}$, zatem k = l, wbrew zalozeniu. Jest wiec, ze $\delta_{k} \neq \delta_{1}$, dla $k \neq l$ i w ten sposob warunek 2. jest spelniony.

Pokazalismy, ze liczby z ciagu (2) co najwyzej porzadkiem roznia sie od liczb ciagu (1).

Wynika stad, ze $\delta_{1}\delta_{2}\dots \delta_{\phi(m)} = r_{1}r_{2}\dots r_{\phi(m)}$

Zatem, wobec (2),

\[ \tag{3} ar_{1}ar_{2}\dots ar_{\phi(m)} = a^{\phi(m)}r_{1}r_{2}\dots r_{\phi(m)} \equiv r_{1}r_{2}\dots r_{\phi(m)} \]  

Niech $P = r_{1}r_{2}\dots r_{\phi(m)}$; poniewaz $gcd(r_{k}, m) = 1$ dla $k = 1,2,\dots, \phi(m)$, wiec bedzie $gcd(P, m) = 1$, a poniewaz wobec (3) mamy $a^{\phi(m)}P \equiv P\text{(mod m)}$, czyli $m|P(a^{\phi(m)}-1)$, wiec musi byc $m| a^{\phi(m)}-1$, czyli $a^{\phi(m)}\equiv 1\text{(mod m)}$, 

Udowodnilismy zatem ze:

\begin{tcolorbox}[colback=white!90!blue,colframe=white!35!blue,title=Twierdzenie Eulera]
    Dla kazdej liczby calkowitej a, ktora jest wzglednie pierwsza do liczby naturalnej m zachodzi kongruencja \[ a^{\phi(m)} \equiv 1\text{(mod m)} \]
 
\end{tcolorbox}

Zauwazmy, ze twierdzenie to jest uogolnieniem malego twierdzenia Fermata, gdyz jesli m jest liczpa pierwsza p, to mamy jak wiadomo, $\phi(m) = \phi(p) = p-1$ a warunek $gcd(a,p) = 1$ jest rownowazny warunkowi $p \nmid a$

Przypomnienie:

\begin{tcolorbox}[colback=white!90!blue,colframe=white!35!blue,title=Male twierdzenie Fermata]

    Jezeli liczba calkowita D jest niepodzielna przez liczbe pierwsza p, to 
    \[ D^{p-1} \equiv 1\text{(mod p)} \]

\end{tcolorbox}

Teraz inne wazne twierdzenie:

\begin{tcolorbox}[colback=white!90!blue,colframe=white!35!blue,title=Funkcja Eulera jest funkcja multiplikatywna]

    Jezeli a i b sa liczbami naturalnymi wzglednie pierwszymi(piszemy $a \perp b$) to:

    \[ \phi(ab) = \phi(a)\phi(b) \tag{4}\]

\end{tcolorbox}

Aby to udowodnic, wykazemy najpierw pomocniczy lemat, zreszta bardzo wazny

\begin{tcolorbox}[colback=white!90!green,colframe=black!35!green,title= Reszty z dzielenia przez a kolejnych liczb od r az do ab - b + r]

    Jezeli a i b sa danymi liczbami naturalnymi takimi, ze $gcd(a,b) = 1$ oraz jezeli r jest dowolna liczba calkowita, to reszty z dzielenia przez a liczb:

    \[r, b+r, b+2r, \dots, (a-1)b + r \tag{5}\]
  
 roznia sie co najwyzej kolejnoscia od liczb:

    \[ 0, 1, 2, \dots, a-1 \tag{6} \]

\end{tcolorbox}

$\textbf{Dowod lematu}$

Zalozmy, ze dla pewnych liczb calkowitych k i l, gdzie $0\leq k < l < a$ liczby kb+r i lb +r daja te sama reszte przy dzieleniu przez a. Byloby wiec $a| lb + r - (kb + r)$ skad $a|(l-k)b$, co wobec $gcd(a, b) = 1$ daje $a| l-k$, co wobe $0\leq k < l < a$ czyli $0 < l-k < a$, jest niemozliwe. Liczby (5) daja wiec same rozne reszty przy dzieleniu przez a, a poniewaz jest ich a, to jest ich tyle ile wszystkich roznych reszt z dzielenia przez a, wiec liczby ciagu (5) roznia sie co najwyzej kolejnoscia od liczb ciagu (6). c.b.d.o.

$\textbf{Dowod twierdzenia o multiplikatywnosci funkcji $\phi$:}$

Wobec $\phi(1) = 1 $ twierdzenie jest prawdziwe, gdy conjamniej jedna z liczb a i b jest rowna 1. Mozemy wiec dalej zakladac, ze a > 1 i b > 1. Liczba $\phi(ab)$ jest liczba tych liczb tablicy(o a wierszach i b kolumnach)

\begin{tabular}{|r|r|r|r|r|r|}
    \hline
    1 & 2 & \dots & r & \dots & b \\
    b+1 & b+2 & \dots & b+r & \dots & 2b \\
    2b+1 & 2b+2 & \dots & 2b+r & \dots& 3b \\
    \vdots & \vdots & \vdots & \vdots & \vdots & \vdots \\
    (a-1)b + 1 & (a-1)b + 2 & \dots &(a-1)b + r & \dots& ab \\
    \hline
\end{tabular}

ktore sa wzglednie pierwsze wzlgedem ab, a wiec ktore sa wzglednie pierwsza z a i wzglednie pierwsze z b.

Niech r bedzie dana liczba naturalna $\leq b$ i wezmy pod uwage r-ta kolumne naszej tablicy. Jezeli $gcd(r,b) = 1$, to wszystkie wyrazy tej kolumny sa wzglednie pierwsze z b, w przeciwnym razie zaden z nich nie jest wzglednie pierwszy z b.

Liczba takich naturalnych $r\leq b$ dla ktorych $gcd(r,b) = 1$  wynosi oczywiscie $\phi(b)$. Tyle tez jest kolumn, ktorych wszystkie wyrazy sa pierwsze wzgledem b. Wezmy teraz ktorakolwiek z takich kolumn, niech to bedzie r-ta kolumna. W mysl naszego lematu(?) reszty z dzielenia przez a liczb wybranej kolumny roznia sie co najwyzej porzadkiem od liczb 0, 1, 2, ..., a-1, skad wynika natychmiast, ze wsrod nich jest $\phi(a)$ liczb pierwszych wzgledem a. W kazdej wiec z $\phi(b)$ kolumn, ktorych wyrazy sa wzglednie pierwsze z b, mamy $\phi(a)$ wyrazow wzglednie pierwszych do a. Wynika stad, ze liczba wyrazow naszej tablicy pierwszych wzgledem a oraz wzgledem b wynosi $\phi(a)\cdot \phi(b)$. Stad wzor (4) i twierdzenie zostalo udowodnione.

Z tego twierdzenia(4) latwo wynika z pomoca indukcji, ze

\begin{tcolorbox}[colback=white!90!blue,colframe=white!35!blue,title=Funkcja $\phi$ jest przeliczalnie multiplikatywna]

    Jezeli $a_{1} a_{2}, \dots, a_{m}$ sa liczbami naturalnymi, parami wzglednie pierwszymi, to

    \[ \phi(a_{1}a_{2}\dots a_{m}) = \phi(a_{1})\phi(a_{2})\dots \phi(a_{m}) \]
\end{tcolorbox}


\begin{problem}{Uogolnione male twierdzenie Fermata}

    Niech $n = p_{1}^{n_{1}}p_{2}^{n_{2}}\dots p_{s}^{n_{s}}$, $\phi$ jest funkcja eulera oraz $\psi(n) = lcm(\phi(p_{1}^{n_{1}}), \phi(p_{2}^{n_{2}}), \dots, \phi(p_{s}^{n_{s}}))$

Udowodnij, ze dla $a\perp n$ zachodzi $n|a^{\psi(n)}-1$

\begin{solution}
    
\end{solution}
\end{problem}

\begin{problem}{Suma funkcji Eulera}
    
    Pokaz, ze $\sum_{d: d|n} \phi(d) = n$

    \begin{solution}
        
    \end{solution}
\end{problem}

\begin{problem}{smh}
    Pokaz wzor: $loglcm(1,2,\dots,n) = \sum_{p \in \mathbb{P}}\lfloor log_p(n) \rfloor log(p)$
    
    \begin{solution}
        
    \end{solution}

\end{problem}

    \begin{problem}{Zasada szufladkowa}
        Udowodnij nastepujace stwierdzenia za pomoca zasady szufladkowej.

        \subproblem
        W turnieju kazda druzyna gra z kazda inna dokladnie raz. W kazdym momencie trwania turnieju istnieja dwie druzyny, ktore rozegraly tyle samo meczow.

        \subproblem

        Wsrod pieciu punktow wybranych w trojkacie rownobocznym o boku 1, istnieje przynajmniej jedna para punktow odleglych od siebie o co najwyzej $\frac{1}{2}$.

        \subproblem

        Kazdy wieloscian wypukly zawiera przynajmniej dwie sciany o tej samej liczbie krawedzi.
        
    \begin{solution}
        

    \end{solution}

    \end{problem}


\end{document}

